\documentclass[11pt]{article}
\usepackage{../shared/hanzocover}
\usepackage[utf8]{inputenc}
\usepackage[T1]{fontenc}
\usepackage{amsmath,amssymb}
\usepackage{graphicx}
\usepackage{booktabs}
\usepackage{xcolor}
\usepackage{listings}
\input{../shared/lstlang}
\usepackage[hidelinks]{hyperref}

\definecolor{codegreen}{rgb}{0,0.6,0}
\definecolor{codegray}{rgb}{0.5,0.5,0.5}
\definecolor{codepurple}{rgb}{0.58,0,0.82}
\definecolor{backcolour}{rgb}{0.96,0.96,0.94}
\definecolor{hzblue}{rgb}{0.12,0.30,0.55}
\definecolor{hzred}{rgb}{0.60,0.16,0.16}

\lstdefinestyle{mystyle}{
    backgroundcolor=\color{backcolour},
    commentstyle=\color{codegreen},
    keywordstyle=\color{codepurple},
    numberstyle=\tiny\color{codegray},
    stringstyle=\color{codegreen},
    basicstyle=\ttfamily\footnotesize,
    breakatwhitespace=false,
    breaklines=true,
    captionpos=b,
    keepspaces=true,
    numbers=none,
    showspaces=false,
    showstringspaces=false,
    showtabs=false,
    tabsize=2,
    frame=single,
    rulecolor=\color{codegray}
}
\lstset{style=mystyle}

\title{The Native Object Store and Its POSIX Projection:\\
       Measured Throughput, Latency, and the Limits of Filesystem Emulation}
\author{Zach Kelling\thanks{Corresponding author: research@hanzo.ai} \\
    \textit{Hanzo Industries Inc (Techstars '17)} \\
    Los Angeles, California \\
    \texttt{https://hanzo.ai}}
\date{Revision~1}

\begin{document}

\hanzocoverpage

\twocolumn[
  \begin{@twocolumnfalse}
  \maketitle
  \begin{abstract}
  \noindent
  We report measurements of \texttt{hanzoai/s3}, the native Hanzo object
  store, and of the POSIX filesystem it projects through a
  Container Storage Interface (CSI) driver whose control path speaks the
  native ZAP transport rather than gRPC. All figures are taken from the
  production deployment, in-cluster, against authenticated requests.

  The object store performs as an object store should: 265.9\,MiB/s on
  16\,MiB reads, 12.81\,ms median metadata lookups, and stable
  latency under mixed concurrent load. The filesystem projection sustains
  287\,MB/s sequential writes and 685\,MB/s sequential reads.

  Its metadata plane does not follow. The mount performs \emph{27 file
  creations per second} and \emph{201 stats per second}, three orders of
  magnitude below local block storage. We quantify the consequence on a real
  metadata-bound workload---a Git repository---and find slowdowns of
  \(27\times\) to \(278\times\). We therefore argue against the common
  instinct to relocate a Git forge onto an object-store-backed mount, and
  state the boundary we believe generalises: \textbf{project an object store
  into POSIX for bulk data, never for namespaces}.

  We also document three defects found while bringing this path to
  production---two in connection-pooled RPC and one silently destroying data
  during metadata restore---because each was invisible to the tests that
  would ordinarily be trusted, and the reasons they were invisible are
  reusable.
  \end{abstract}
  \vspace{1em}
  \end{@twocolumnfalse}
]

\section{System Under Test}

\texttt{hanzoai/s3} is a fork of the SeaweedFS lineage that has replaced its
filer RPC surface with ZAP, the Hanzo capability transport: 28 unary and 5
streaming methods answer over \texttt{transport.ListenStream} on the port
historically named the ``gRPC port''. No protobuf framing remains on that
path. The measured build is \texttt{v1.0.4}.

The metadata store is \texttt{luxdb}, a wrapper over
\texttt{github.com/luxfi/database} with a \texttt{zapdb} backend. This is the
sole store compiled into the binary; the SeaweedFS \texttt{leveldb2} store was
removed. That removal is consequential and we return to it in
Section~\ref{sec:migration}.

The POSIX projection is \texttt{hanzoai/s3-csi}, a CSI driver forked so that
its filer-facing half speaks ZAP. A volume is a directory beneath the filer
root, so volumes are \texttt{ReadWriteMany} by construction. The
kubelet-facing half remains CSI gRPC, which is the Kubernetes contract rather
than a design choice.

\subsection{Deployment Shape, Stated Honestly}
\label{sec:durability}

It is worth being exact about durability, because the architecture is often
described more ambitiously than the deployment warrants.

\begin{itemize}
  \item Volume replication is configured \texttt{"000"}: \textbf{one copy of
        every volume}. There is no redundancy at the store layer.
  \item The topology contains \textbf{one} volume server and one master
        member.
  \item The Deployment is \texttt{Recreate}, \texttt{replicas: 1}, bound to a
        single 700\,GiB \texttt{ReadWriteOnce} block volume.
  \item \texttt{luxfi/consensus} is \textbf{not} a dependency of this
        service. No consensus protocol participates in its durability.
\end{itemize}

Durability therefore rests entirely on the underlying cloud block device and
its snapshots. The store is not highly available: a roll is an outage, and a
lost volume is a restore-from-backup event. We state this plainly because a
paper that implied otherwise would be describing a system we do not run.
Introducing replication, or placing the metadata plane under a consensus
protocol, remains future work rather than shipped behaviour.

\section{Method}

Measurements were taken from pods inside the same cluster as the store, to
exclude ingress and WAN effects. Object-store figures use SigV4-signed
requests through the service DNS name; the driver and the store were the
production instances serving live traffic. Each latency figure is the mean of
ten repetitions; metadata figures are means over 20--500 operations.
Filesystem figures use \texttt{dd} with \texttt{conv=fsync} for writes.

The Git comparison runs an identical script against the CSI mount and against
the pod's local disk, so the delta isolates the storage path.

\section{Object Store Results}

\begin{table}[h]
\centering
\begin{tabular}{lrrr}
\toprule
Object size & PUT (ms) & GET (ms) & GET (MiB/s) \\
\midrule
1\,KiB   &  38.93 & 15.67 &   0.1 \\
64\,KiB  &  18.21 & 17.10 &   3.7 \\
1\,MiB   &  36.46 & 20.30 &  49.3 \\
16\,MiB  & 244.80 & 60.17 & 265.9 \\
\bottomrule
\end{tabular}
\caption{Signed request latency, in-cluster, mean of 10.}
\end{table}

Metadata lookups (\texttt{HEAD} on an existing object) average
12.81\,ms over 20 repetitions. Latency is dominated by a fixed
per-request cost of roughly 15\,ms until objects exceed 1\,MiB,
after which transfer dominates and throughput climbs to 265.9\,MiB/s.

Under a sustained mixed load---eight rounds of four concurrent
1000-key \texttt{LIST} operations interleaved with
\texttt{PUT}/\texttt{GET}/\texttt{DELETE}---the store completed every
operation with no failures and unary latency flat between 7.5 and
21\,ms. This load shape is precisely what exposed the transport defect
of Section~\ref{sec:hol}, and its stability is the evidence that the defect
is resolved.

\section{Filesystem Projection Results}

\begin{table}[h]
\centering
\begin{tabular}{lrr}
\toprule
Operation & Block size & Rate \\
\midrule
Sequential write & 1\,MiB & 188\,MB/s \\
Sequential write & 4\,MiB & 287\,MB/s \\
Sequential read  & 1\,MiB & 651\,MB/s \\
Sequential read  & 4\,MiB & 685\,MB/s \\
\midrule
File creation & --- & \textbf{27 files/s} \\
\texttt{stat} & --- & \textbf{201 stats/s} \\
\bottomrule
\end{tabular}
\caption{CSI mount, 20\,GiB RWX volume.}
\end{table}

The bulk-data path is strong and would satisfy most streaming workloads. The
namespace path is not: creating a file costs approximately 37\,ms and
\texttt{stat} approximately 5\,ms. Both are network round trips to the
filer with no client-side namespace cache that survives them.

\section{The Metadata Cost, Made Concrete}

Aggregate rates understate how punishing this is, because real workloads
issue metadata operations in dependent chains rather than in parallel. We
therefore measured Git, which is metadata-bound almost by definition: loose
objects, refs, and index updates are small files created and stat'ed in
sequence.

\begin{table}[h]
\centering
\begin{tabular}{lrrr}
\toprule
Operation & Local & CSI mount & Slowdown \\
\midrule
init + 200 files + commit & 1617\,ms & 43496\,ms & \(27\times\) \\
\texttt{clone}            & 83\,ms   & 23044\,ms & \(278\times\) \\
\texttt{gc --aggressive}  & 162\,ms  & 9699\,ms  & \(60\times\) \\
\texttt{log} + \texttt{status} & 12\,ms & 2857\,ms & \(238\times\) \\
\bottomrule
\end{tabular}
\caption{Identical Git workload, same pod, two storage paths.}
\end{table}

A 278-fold penalty on \texttt{clone} is not a tuning problem. It is
the arithmetic of thousands of dependent round trips at 5--37\,ms
apiece.

\subsection{Consequence for the Git Forge}

The motivating goal of this work was to relocate a 22\,GiB Git forge
(329 bare repositories) off a single \texttt{ReadWriteOnce} volume---the
constraint pinning the control-plane binary to one node---onto shared
storage. The forge's storage layer is an \texttt{osfs} rooted at a data
directory, so the change appeared to require no code at all: mount the
directory elsewhere.

The measurements say do not. Git's heavy paths (clone/fetch serve, push
receive) shell out to the \texttt{git} CLI against an OS path and stream
multi-gigabyte packfiles, which this mount serves well. But every other Git
operation is metadata, and metadata is where the projection collapses.

The honest conclusions are:
\begin{enumerate}
  \item \textbf{Do not place the live Git forge on this mount.} The
        \(238\times\) penalty on \texttt{log}/\texttt{status} alone would be
        user-visible on every request.
  \item \textbf{The object store is the right home for packfiles}, which are
        large, sequential, and immutable---exactly the shape it serves at
        265.9\,MiB/s.
  \item \textbf{Decouple the two problems.} Single-node pinning is a
        \emph{placement} problem; solving it with a filesystem projection
        imports a metadata cost the workload cannot absorb. A Git-native
        object backend, or replication at the block layer, addresses
        placement without paying it.
\end{enumerate}

We regard this as the paper's main result, and note it inverts the plan we
began with. The prior reasoning---that the code was already correct and only
the mount was missing---was true about the code and wrong about the
workload.

\section{Defects Found En Route}

Three defects had to be fixed before any of the above could be measured.
Each is reported because each defeated the check that should have caught it.

\subsection{Promise identifiers scoped to the client, not the connection}

Generated ZAP clients each carried a private \texttt{rpc.Session} numbering
calls from~1, while the connection pool shares one connection per peer
address and constructs a fresh client per call. Two concurrent calls
therefore both claimed promise~1; the second registration evicted the first
from the connection's in-flight table, and the evicted response was discarded
on arrival. The caller blocked for the lifetime of its context---at service
start-up, \texttt{context.Background()}.

The invariant: \textbf{a promise identifier is the connection's
demultiplexing key, so the connection must allocate it}. Stream identifiers
already did.

\subsection{Metadata restore that dropped every file}

The filer's \texttt{meta.load} created directories inline and files on
goroutines, and returned on end-of-input without waiting for them. The caller
closes the connection on return, so in-flight writes died against a torn-down
connection---and the command reported success. A restored tree came back with
its directory structure intact and \emph{no files}.

This is the failure mode one discovers during a disaster recovery. It
survived because the only signal was the absence of data nobody had yet
tried to read.

\subsection{Head-of-line blocking on pooled connections}
\label{sec:hol}

Inbound stream frames were delivered from the connection's read loop with a
blocking send into a 16-slot buffer. A consumer slower than its producer
filled that buffer, stalled the read loop, and with it every other stream and
every unary reply on the connection. Pooling turned a throughput concern into
an outage: one large \texttt{LIST} froze every request in the process.

It presented as an authentication fault, because unauthenticated requests are
rejected before reaching the filer and continued answering in
5\,ms. Diagnosis required a goroutine dump from the stalled process; no
log line existed, because the code was not failing, only waiting.

The fix queues per stream and wakes the reader, so delivery never blocks the
read loop; a bound fails one stream rather than the connection everyone
shares.

\subsection{What the three have in common}

Two of the three were latent properties of \emph{connection pooling}---a
sharing optimisation that converts an isolated stall into a global one---and
all three were invisible to unauthenticated smoke tests, to fresh-dataset
tests, and to green unit suites. The tests that found them were: an
adversarial reproduction with a deliberately slow consumer; a restore
exercised against data written by the \emph{previous} release; and a
goroutine dump taken while production was hung.

\section{Store Migration}
\label{sec:migration}

Because \texttt{luxdb} is the only store compiled in, a version upgrade
against an existing \texttt{leveldb2} data directory starts with an
\textbf{empty namespace}---no error, no crash, every bucket invisible. The
two layouts coexist in one directory (\texttt{leveldb2} in numbered
sub-directories, \texttt{zapdb} in top-level files), which is what makes the
migration reversible.

The production migration moved 50{,}006 directories and 68{,}646 files, and
all 45 buckets returned. Reversal is an image change; only writes made in the
new store are lost. The upgrade path is therefore: export, roll, import,
verify with a \emph{signed} operation against \emph{pre-existing} data---the
last qualifier being the one that matters, since unauthenticated and
freshly-created-object checks pass while the store is broken.

\section{Conclusion}

The native object store is fast, and its POSIX projection is fast for bulk
data and slow for namespaces by three orders of magnitude. That asymmetry is
not an implementation defect to be optimised away; it is the cost of
emulating a filesystem over a network metadata service, and it decides which
workloads may cross the boundary.

Bulk, sequential, immutable data belongs in the object store. Namespaces with
dependent metadata chains---Git foremost---do not, and should not be moved
there merely because their storage layer makes it syntactically easy.

\end{document}
